\documentclass[12pt,a4paper]{article}
\usepackage{fontspec}
\setmainfont{Times New Roman}
\usepackage[top=2cm,bottom=2cm,left=2.5cm,right=2cm]{geometry}
\usepackage{enumitem}
\usepackage{tabularx}
\usepackage{booktabs}
\usepackage{array}
\newcolumntype{Y}{>{\centering\arraybackslash}X}
\pagestyle{plain}
\linespread{1.15}

\begin{document}
% --- HEADER 2 CỘT CHUẨN NGHỊ ĐỊNH 30/2020/NĐ-CP ---
\noindent
\begin{minipage}[t]{0.45\textwidth}
    \begin{center}
        \fontsize{11pt}{13pt}\selectfont
        CÔNG AN THÀNH PHỐ HÀ NỘI\\
        \textbf{PHÒNG CẢNH SÁT HÌNH SỰ (PC02)}\\[-0.2em]
        \rule{3.2cm}{0.6pt}\\[0.25cm]
        \fontsize{11pt}{13pt}\selectfont Số: \textbf{02/BB-HC}
    \end{center}
\end{minipage}
\hfill
\begin{minipage}[t]{0.52\textwidth}
    \begin{center}
        \fontsize{11pt}{13pt}\selectfont
        \textbf{CỘNG HÒA XÃ HỘI CHỦ NGHĨA VIỆT NAM}\\
        \textbf{Độc lập - Tự do - Hạnh phúc}\\[-0.2em]
        \rule{3.6cm}{0.6pt}\\[0.25cm]
        \textit{Hà Nội, ngày 25 tháng 07 năm 2016}
    \end{center}
\end{minipage}

\vspace{0.6cm}

\begin{center}
    \textbf{\Large BIÊN BẢN LẤY LỜI KHAI}\\[0.15cm]
    \textit{}
\end{center}

\vspace{0.4cm}
\textit{(Lần 2 đối với đối tượng Lê Quang Vũ)}



Vào hồi 17 giờ 45 phút, ngày 25 tháng 07 năm 2016, tại Phòng Cảnh sát Hình sự Công an TP. Hà Nội.

\textbf{Chúng tôi gồm:}

\begin{enumerate}[label=\arabic*., leftmargin=1.2cm, itemsep=2pt]
  \item Điều tra viên: Đại úy Lê Minh — Đội Điều tra Trọng án.
  \item Cán bộ ghi biên bản: Trung úy Nguyễn Văn Hoàng — Cán bộ Đội Án mạng.

\end{enumerate}
\textbf{Tiến hành lấy lời khai lần 02 đối tượng:}

\begin{itemize}[leftmargin=1.2cm, itemsep=2pt]
  \item \textbf{Họ và tên:} \textbf{LÊ QUANG VŨ} | \textbf{Giới tính:} Nam | \textbf{Sinh ngày:} 20/09/1991.
  \item \textbf{Nơi ở hiện nay:} Số 45, Đường Đoàn Kết, Phường Cảng Đông, TP. Hà Nội.


\end{itemize}

\subsection*{NỘI DUNG HỎI CUNG (HỎI VÀ ĐÁP)}


\noindent\textbf{Hỏi (ĐTV Lê Minh):} \textit{Chúng tôi đã có trong tay dữ liệu:}\\[0.15cm]
\begin{enumerate}[label=\arabic*., leftmargin=1.2cm, itemsep=2pt]
  \item Lời khai nhân chứng xác nhận vợ anh (Nguyễn Ngọc Mai) nổ máy xe rời đi lúc \textbf{19:00}.
  \item Ảnh chụp màn hình ứng dụng xe công nghệ trên máy của anh ghi nhận thời gian gửi lệnh đặt xe là \textbf{19:25:40} và tài xế đón lúc \textbf{19:30:15}.
  \item Cuốn sổ nợ của Khang ghi rõ con nợ \texttt{Thằng Lệch Pha} nợ \texttt{300.000.000đ} (SĐT \texttt{0967.452.183} — chúng tôi đã xác minh đây là số điện thoại thứ hai đăng ký dưới tên anh, đồng thời chữ ký trong sổ nợ tương đồng so với chữ ký của anh trong biên bản lấy lời khai đầu tiên; và biệt danh \texttt{Lệch Pha} trong sổ nợ trùng khớp với biệt danh Khang gọi anh trong các tin nhắn SMS gửi đến số \texttt{0988.20.09.91}).
\end{enumerate}
Dựa vào những chứng cứ trên, Cơ quan điều tra có đủ cơ sở xác định anh đã ở lại nhà anh Khang thêm 30 phút sau khi vợ anh rời đi.

Anh giải thích thế nào về việc này và vì sao anh lại giấu chi tiết này trong lời khai đầu tiên?


\noindent\textbf{Đáp (Lê Quang Vũ):} \textit{(Mặt tái mét, ôm đầu bật khóc nức nở)} Dạ... thưa cán bộ... tôi xin khai thật hết... Tôi nhận thầu điện bị đọng vốn nên có giấu vợ vay nặng lãi của anh Khang 300 triệu từ đầu năm. Hạn trả là ngày 20/07 vừa rồi nhưng tôi không xoay được tiền. Anh Khang tức giận nhắn tin đe dọa sẽ ném giấy nợ về cho bố mẹ vợ tôi và cơ quan vợ tôi biết. Tôi ở rể, sợ mất mặt và sợ vợ ly hôn nên tôi hoảng loạn tột cùng...\\[0.35cm]

\noindent\textbf{Hỏi (ĐTV Lê Minh):} \textit{Tối 24/07, trong 30 phút ở lại từ 19:00 đến 19:30, anh đã làm gì trong nhà nạn nhân Khang?}\\[0.15cm]

\noindent\textbf{Đáp (Lê Quang Vũ):} Lúc vợ tôi tức giận ném giấy tờ rồi nổ máy xe về lúc 19:00, tôi cố tình nói dối Mai là đi nhậu với bạn để Mai về trước. Sau khi Mai về , tôi vội đóng cửa bước vào phòng khách quỳ lạy xin anh Khang cho tôi hoãn nợ. Nhưng anh Khang cười khẩy, rút cuốn sổ da đen đập mạnh vào ngực tôi chửi: \textit{"Mày bùng hẹn ngày 20 rồi, hạn chót ngày mai không nôn đủ 300 củ tao in 100 tờ rơi rải khắp phố Đoàn Kết cho bố vợ mày nhục mặt!"}.\\[0.35cm]
Không xin gia hạn được món nợ, tôi cũng chẳng buồn nán lại thêm làm gì. Vì tâm trạng không tốt, tôi quyết định bắt xe đi nhậu một mình... À, lúc đứng ở đầu ngõ đợi xe, tôi có để ý thấy một bóng người mặc áo gió trùm mũ kín mít đang đứng nép dưới gốc cây ngó vào nhà anh Khang... Nhưng lúc đó xe đã đến nên tôi cũng lên xe luôn mà không để ý gì thêm...

\noindent\textbf{Hỏi (ĐTV Lê Minh):} \textit{Anh có gì để chứng minh mình ngồi nhậu đến hết tối mà không phải tranh thủ cơ hội quay lại ra tay với nạn nhân?}\\[0.15cm]
\noindent\textbf{Đáp (Lê Quang Vũ):} Tôi bắt xe ôm ra Quán Bia 88 ở Cầu Cảng ngồi một mình ở góc sát mép sông, gọi 1 đĩa nem chua với 5 lon bia. Buồn đời nên nốc bia liên tục say bí bét chẳng để ý giờ giấc gì cả... Hình như tầm tám rưỡi gì đấy là tôi đứng dậy thanh toán thì phải, tôi chả nhớ rõ nữa... Lúc về tôi chỉ nhớ còn thấy ngay trước cửa quán có vụ xô xát giữa mấy thanh niên, tôi đứng ngó một lúc rồi mới bắt xe ôm về. Tôi thề là tôi về thẳng nhà tôi thôi chứ không quay lại nhà anh Khang.\\[0.35cm]

\noindent\textbf{Hỏi (ĐTV Lê Minh):} \textit{Vậy lúc anh rời khỏi nhà lúc 19:25, tình trạng nạn nhân Khang thế nào?}\\[0.15cm]
\noindent\textbf{Đáp (Lê Quang Vũ):} Lúc tôi bỏ đi, anh Khang vẫn còn đứng ở giữa phòng khách chửi bới ầm ĩ và đá cái ghế gỗ. Tôi không hề giết anh Khang, dù sao đi nữa anh ấy cũng là anh vợ của tôi mà cán bộ...\\[0.35cm]

À... mà thưa cán bộ... nếu... nếu anh Khang chết rồi thì... khoản nợ 300 triệu đó tôi có phải trả nữa không ạ?



Biên bản hỏi cung kết thúc vào hồi 18 giờ 30 phút cùng ngày. Biên bản đã được đọc lại cho Lê Quang Vũ nghe, công nhận đúng sự thật và cùng ký tên.


\textbf{ĐIỀU TRA VIÊN}

\textit{(Ký tên)}

\textbf{Đại úy Lê Minh}


\textbf{NGƯỜI BỊ HỎI CUNG}

\textit{(Ký, ghi rõ họ tên)}

\textbf{Lê Quang Vũ}

\end{document}
